% Бүлэг 3

\chapter{Системийн зохиомж ба хэрэгжүүлэлт}
\label{Chapter3}
\pagecolor{white}

%-------------------------------------------------------------------------------

\section{Системийн архитектур}
\textbf{FinTrack} систем нь сонгодог \textbf{3-tier} архитектурын дээр AI давхаргыг нэмж, дөрвөн давхаргат архитектураар бүтсэн \cite{cleanarch, ddd}. Энэ нь дараах ач холбогдолтой:
\begin{itemize}
    \item \textbf{Хариуцлагын тусгаарлалт.} Дэлгэц, бизнесийн логик, өгөгдөл хадгалалт, AI тус бүр өөрийн модультай.
    \item \textbf{Засвар үйлчилгээ хялбар.} Нэг давхаргыг өөрчилсөн нь нөгөөд бараг нөлөөлөхгүй.
    \item \textbf{Тестлэх боломжтой.} Давхарга бүрд нэгж тест бичиж болно.
    \item \textbf{Өргөтгөх боломж.} Сервер талыг хувилж олон нэгж дээр зэрэгцэн ажиллуулах боломжтой.
\end{itemize}

\begin{figure}[ht]
\centering
\begin{tikzpicture}[
    box/.style={rectangle, rounded corners, draw=darknavy, very thick, minimum width=12cm, minimum height=1.4cm, align=center, font=\footnotesize},
    presentation/.style={box, fill=accentblue!20},
    application/.style={box, fill=accentgreen!20},
    data/.style={box, fill=accentorange!20, minimum width=5.7cm},
    aiext/.style={box, fill=accentpurple!20, minimum width=5.7cm},
    arrow/.style={->, very thick, draw=darknavy}
]
\node[presentation] (p) at (0, 6) {%
    \textbf{Танилцуулах давхарга (Presentation)} \\
    \scriptsize React Native + Expo SDK 55 + TypeScript + NativeWind v4 \\
    \scriptsize 17 дэлгэц $\cdot$ React Navigation $\cdot$ AuthContext
};
\node[application] (a) at (0, 3.5) {%
    \textbf{Бизнес логикийн давхарга (Application)} \\
    \scriptsize Go 1.25 + Gin + JWT дундын давхарга $\cdot$ 12 боловсруулагч файл $\cdot$ 45 REST API цэг \\
    \scriptsize Үйлдлийн бүртгэл $\cdot$ Expo push дамжуулагч $\cdot$ Өдөр тутмын AI хуваарьчлагч
};
\node[data] (d) at (-3.2, 1) {%
    \textbf{Өгөгдлийн давхарга (Data)} \\
    \scriptsize PostgreSQL 16 (port 5433) $\cdot$ GORM v1.25.9 $\cdot$ 11 хүснэгт \\
    \scriptsize \textbf{MinIO} S3-нийцтэй обьект хадгалалт (хуулга, аватар)
};
\node[aiext] (ext) at (3.2, 1) {%
    \textbf{AI / гадаад үйлчилгээний давхарга} \\
    \scriptsize Gemini 2.5 / OpenRouter (нөөц) \\
    \scriptsize FastAPI хуулга задлагч $\cdot$ Expo Push
};

\draw[arrow] (p) -- node[right, font=\scriptsize]{HTTPS REST + JWT} (a);
\draw[arrow] (a.south) -- ++(0, -0.7) -| (d.north);
\draw[arrow] (a.south) -- ++(0, -0.7) -| (ext.north);
\end{tikzpicture}
\caption{4-давхаргат архитектур}
\label{fig:arch-layers}
\end{figure}

\subsection{Сервер талын кодын зохион байгуулалт}
Сервер тал нь Go-ын давхаргат зохион байгуулалтаар (\texttt{internal/} pattern) бүтсэн \cite{cleancode, dpo}:

\begin{table}[ht]
\centering
\footnotesize
\renewcommand{\arraystretch}{1.25}
\rowcolors{2}{softgrey}{white}
\begin{tabular}{|p{4.2cm}|p{8cm}|}
\hline
\rowcolor{darknavy}
\textcolor{white}{\textbf{Хавтас}} & \textcolor{white}{\textbf{Зориулалт}} \\ \hline
\texttt{cmd/server/main.go}        & Гарааны цэг. Тохиргоо ачаалах, өгөгдлийн санд холбогдох, чиглүүлэлт бүртгэх, хуваарьчлагч эхлүүлэх. \\ \hline
\texttt{internal/config}           & Орчны хувьсагч уншигч (\texttt{godotenv}). \\ \hline
\texttt{internal/database}         & PostgreSQL холболт, автомат миграц, эх өгөгдөл. \\ \hline
\texttt{internal/middleware}       & JWT баталгаажуулалт, CORS. \\ \hline
\texttt{internal/handlers}         & 11 файл --- боловсруулагчийн логик. \\ \hline
\texttt{internal/models}           & 9 GORM загвар + DTO. \\ \hline
\texttt{internal/routes}           & Чиглүүлэлт бүртгэгч. \\ \hline
\end{tabular}
\caption{Сервер талын кодын зохион байгуулалт}
\label{tab:backend-layout}
\end{table}

\subsection{Боловсруулагчдын жагсаалт}

\begin{table}[ht]
\centering
\footnotesize
\renewcommand{\arraystretch}{1.25}
\rowcolors{2}{softgrey}{white}
\begin{tabular}{|p{3.5cm}|c|p{6.5cm}|}
\hline
\rowcolor{darknavy}
\textcolor{white}{\textbf{Файл}} & \textcolor{white}{\textbf{Мөр}} & \textcolor{white}{\textbf{Үндсэн зорилго}} \\ \hline
\texttt{auth.go}            & 316  & Бүртгэл, нэвтрэлт, JWT, профайл, нийгмийн сүлжээний нэвтрэлт. \\ \hline
\texttt{transaction.go}     & 223  & Гүйлгээ CRUD, үлдэгдлийн автомат шинэчлэлт. \\ \hline
\texttt{dashboard.go}       & 288  & Хяналтын самбарын үзүүлэлт, зарлагын нэгтгэл, статистик. \\ \hline
\texttt{budget.go}          & 138  & Сарын төсөв CRUD. \\ \hline
\texttt{account.go}         & 184  & Данс CRUD, ангиллын CRUD. \\ \hline
\texttt{ai\_chat.go}        & 548  & AI чат сесс, мессеж, LLM дуудалт, контекст бэлтгэл. \\ \hline
\texttt{ai\_provider.go}    & 388  & Олон нийлүүлэгчийн солилт (Gemini $\leftrightarrow$ OpenRouter), нөөц шилжилтийн логик, үнэгүй загварын жагсаалт автомат татах. \\ \hline
\texttt{ai\_analysis.go}    & 1{,}425 & Хуулга байршуулах, задлан унших, AI хураангуй гаргах. Хамгийн том боловсруулагч. \\ \hline
\texttt{import.go}          & 197  & Хуучин текст-only хувилбар. \\ \hline
\texttt{notification.go}    & 279  & Push мэдэгдлийн токен, өдөр тутмын автомат хуваарьчлагч. \\ \hline
\texttt{scheduled\_payment.go} & 103 & Давтагдах төлбөр CRUD, идэвх солих. \\ \hline
\texttt{activity\_log.go}   & 88   & Аудит лог жагсаалт ба нэгтгэл. \\ \hline
\end{tabular}
\caption{Сервер талын боловсруулагчдын тойм (нийт $\approx$4{,}900 мөр Go код)}
\label{tab:handlers}
\end{table}

\subsection{Үндсэн загвар ба харилцааны (Class) диаграм}
Сервер талын \texttt{internal/models} пакет дахь үндсэн ангийн (struct) бүтэц, тэдгээрийн талбар, метод болон хоорондын харилцааг доорх ангийн диаграмаар үзүүлэв. Энэ нь ER загвартай ойролцоо боловч програм хангамжийн түвшинд буюу GORM-ийн загварын давхаргыг тусгасан.

\begin{figure}[ht]
\centering
\resizebox{\textwidth}{!}{%
\begin{tikzpicture}[
    >={Stealth[length=2.2mm]},
    class/.style={draw=darknavy, fill=accentblue!15, very thick, rectangle, align=left, inner sep=4pt, font=\scriptsize\ttfamily, text width=3.7cm},
    classfin/.style={class, fill=accentgreen!18},
    classai/.style={class, fill=accentpurple!18},
    classaud/.style={class, fill=accentorange!22},
    assoc/.style={->, semithick, draw=darknavy!85},
    compose/.style={{Diamond[length=3.2mm,width=2.4mm]}-{Stealth[length=2.2mm]}, semithick, draw=darknavy!85},
    card/.style={font=\tiny, fill=white, inner sep=1pt, rounded corners=1pt}
]

% =================== Row 1: Category (зүүн дээд), User (төв) ===================
\node[class] (cat) at (-6, 0) {%
\textbf{Category} \\
\hrulefill \\
- ID : uint \\
- Name, Icon, Color : string \\
- Type : string \\
\hrulefill \\
+ IsIncome() : bool};

\node[class] (user) at (10, 0) {%
\textbf{User} \\
\hrulefill \\
- ID : uint \\
- Name, Email : string \\
- Password : string \\
- Currency, PushToken : string \\
\hrulefill \\
+ HashPassword() \\
+ CheckPassword() : bool \\
+ GenerateJWT() : string};

% =================== Row 2: 6 хэрэглэгчийн хүүхэд анги ===================
\node[classfin] (acc)   at (-2.5, -6) {%
\textbf{Account} \\
\hrulefill \\
- ID, UserID : uint \\
- Name, Type, Icon : string \\
- Balance : float64 \\
\hrulefill \\
+ UpdateBalance(amt) \\
+ Validate() : error};

\node[classfin] (bud)   at (2.5, -6) {%
\textbf{Budget} \\
\hrulefill \\
- ID, UserID : uint \\
- CategoryID : uint \\
- Amount, Spent : float64 \\
- Month, Year : int \\
\hrulefill \\
+ Remaining() : float64 \\
+ IsExceeded() : bool};

\node[classai] (chat)   at (7.5, -6) {%
\textbf{AIChat} \\
\hrulefill \\
- ID, UserID : uint \\
- Title : string \\
- CreatedAt : time.Time \\
\hrulefill \\
+ AddMessage(msg)};

\node[classai] (anal)   at (12.5, -6) {%
\textbf{AIAnalysis} \\
\hrulefill \\
- ID, UserID : uint \\
- Filename, BankName : string \\
- Income, Expenses : float64 \\
- CategoriesJSON : json \\
\hrulefill \\
+ Summary() : string};

\node[classaud] (notif) at (17.5, -6) {%
\textbf{Notification} \\
\hrulefill \\
- ID, UserID : uint \\
- Title, Body, Type : string \\
- IsRead : bool \\
\hrulefill \\
+ MarkRead()};

\node[classaud] (log)   at (22.5, -6) {%
\textbf{ActivityLog} \\
\hrulefill \\
- ID, UserID : uint \\
- Action, Entity : string \\
- Details, IPAddress : string \\
- Status : string};

% =================== Row 3: ач хүүхэд анги (зэрэгцүүлсэн) ===================
\node[classfin] (tx)    at (-2.5, -12) {%
\textbf{Transaction} \\
\hrulefill \\
- ID, UserID : uint \\
- AccountID, CategoryID : uint \\
- Amount : float64 \\
- Date : time.Time \\
- Description : string \\
\hrulefill \\
+ ApplyToBalance() \\
+ Reverse()};

\node[classfin] (sch)   at (2.5, -12) {%
\textbf{ScheduledPayment} \\
\hrulefill \\
- ID, UserID, AccountID : uint \\
- CategoryID : uint \\
- Frequency : string \\
- NextDate : time.Time \\
- IsActive : bool \\
\hrulefill \\
+ Trigger() \\
+ Toggle()};

\node[classai] (msg)    at (7.5, -12) {%
\textbf{AIMessage} \\
\hrulefill \\
- ID, ChatID : uint \\
- Role : string \\
- Content : text \\
\hrulefill \\
+ IsUser() : bool};

% =================== Холболтууд ===================

% --- User-ийн компози харилцаа (User → 6 хүүхэд анги) ---
\draw[compose] (user.south) -- node[card, pos=0.55]{1..*} (acc.north);
\draw[compose] (user.south) -- node[card, pos=0.6]{1..*}  (bud.north);
\draw[compose] (user.south) -- node[card, pos=0.65]{1..*} (chat.north);
\draw[compose] (user.south) -- node[card, pos=0.65]{1..*} (anal.north);
\draw[compose] (user.south) -- node[card, pos=0.6]{1..*}  (notif.north);
\draw[compose] (user.south) -- node[card, pos=0.55]{1..*} (log.north);

% --- Account-ийн компози харилцаа (босоо зэрэгцүүлэлттэй) ---
\draw[compose] (acc.south) -- node[card]{1..*} (tx.north);
\draw[compose] (acc.south east) -- ([xshift=-0.8cm]sch.north);

% --- AIChat → AIMessage (босоо шууд) ---
\draw[compose] (chat.south) -- node[card]{1..*} (msg.north);

% --- Category-ийн ассоциаци (орт-маршрутаар цэгцлэв) ---
% Cat → Budget: дээгүүр гарч Budget-ийн оройд бууна
\draw[assoc] (cat.east) -- (2.5, 0) -- node[card, pos=0.95]{1..*} (bud.north);
% Cat → Transaction: зүүн талын коридороор бууж зүүн дээд буланд бууна
\draw[assoc] (cat.south) -- (-6, -10.5) -- node[card, pos=0.6]{1..*} ([xshift=-0.3cm]tx.north);
% Cat → ScheduledPayment: Budget ба AIChat-ийн хоорондох зайгаар буцаж бууна
\draw[assoc] (cat.east) -- (5, 0) -- (5, -10.5) -- node[card, pos=0.5]{1..*} ([xshift=0.3cm]sch.north);

\end{tikzpicture}
}
\caption{Class Diagram (Класс диаграм) --- сервер талын загварууд ба харилцаа (композици: $\blacklozenge\!-$, ассоциаци: $\rightarrow$)}
\label{fig:class}
\end{figure}

\textbf{Тэмдэглэгээ.} Ромб ($\blacklozenge$) тэмдэгтэй холболт нь \textbf{композици} харилцаа --- эзэмшигч (User эсвэл Account) устгахад дэд объектууд хамт устдаг (ON DELETE CASCADE-той тохирно). Сум ($\rightarrow$) нь Category-аас Transaction, Budget, ScheduledPayment рүү чиглэх \textit{ассоциаци} (лавлагаа) харилцааг илэрхийлнэ. Хэрэглэгч бүх модулийн төв байх тул дээд эгнээнд төвлөрсөн, доошоо зургаан хүүхэд анги салаалсан хэлбэрээр байрласан.

\section{Өгөгдлийн сангийн зохиомж}
Систем PostgreSQL 16-г сонгосон бөгөөд 11 хүснэгтэд автомат миграц хийгдэнэ.\\

PostgreSQL-ийг сонгосон гол шалтгаан \cite{kleppmann}:
\begin{itemize}
    \item \textbf{ACID шинж.} Санхүүгийн системд транзакцийн бүрэн бүтэн байдал чухал.
    \item \textbf{Foreign key хатуу хэрэгжүүлэлт.} Нэг хэрэглэгчийн гүйлгээ нь зөвхөн өөрийнх нь дансанд харьяалагдана.
    \item \textbf{JSON багана.} Уян хатан бүтэцтэй өгөгдлийг (ангиллын задаргаа JSON) шууд хадгалах боломжтой.
    \item \textbf{Индексүүд.} Хайлт болон хуудаслалт хурдан байх.
    \item \textbf{Бодит зэрэгцээ хандалтын хяналт.} Олон хэрэглэгчээс зэрэг хандсан тохиолдолд тохирох түгжих (locking) механизмтай.
\end{itemize}

\subsection{Хүснэгт бүрийн зориулалт}
\begin{table}[ht]
\centering
\footnotesize
\renewcommand{\arraystretch}{1.25}
\rowcolors{2}{softgrey}{white}
\begin{tabular}{|l|p{3cm}|p{6.5cm}|}
\hline
\rowcolor{darknavy}
\textcolor{white}{\textbf{Хүснэгт}} & \textcolor{white}{\textbf{Цөм багана}} & \textcolor{white}{\textbf{Зориулалт}} \\ \hline
\texttt{users}        & email, password         & Хэрэглэгчийн данс. \\ \hline
\texttt{accounts}     & balance, type           & Олон данс (4 төрөл). \\ \hline
\texttt{categories}   & name, type              & Орлого/зарлагын ангилал (17 эх өгөгдөл). \\ \hline
\texttt{transactions} & amount, date            & Бүх гүйлгээ. \\ \hline
\texttt{budgets}      & month, year, amount     & Сар тутмын төсөв ангиллаар. \\ \hline
\texttt{scheduled\_payments} & frequency        & Давтагдах төлбөр. \\ \hline
\texttt{ai\_chats}    & title                   & Чат сесс. \\ \hline
\texttt{ai\_messages} & role, content          & Чат мессеж (user/assistant). \\ \hline
\texttt{ai\_analyses} & categories\_json        & Банкны хуулгын шинжилгээ. \\ \hline
\texttt{notifications} & title, body, type      & Push мэдэгдлүүд. \\ \hline
\texttt{activity\_logs} & action, ip\_address  & Аудитын мөр. \\ \hline
\end{tabular}
\caption{Өгөгдлийн санд хадгалагдах 11 хүснэгт}
\label{tab:tables}
\end{table}

\subsection{Хүснэгтүүдийн уялдаа (ER диаграм)}
ER диаграммыг харьцангуй цэгцтэй байлгахын тулд хүснэгтүүдийг 4 эгнээ, 3 баганад байрлуулж, гадаад түлхүүрийн уялдааг \textbf{зөв өнцөгт сум} буюу зөв өнцгөөр эргэх шугамаар зурсан. Уялдааны төгсгөлд тэмдэг тавиагүй учир бүх холболт нь \textit{1:N} (нэг талд олон) харьцаатай.

\begin{figure}[ht]
\centering
\resizebox{\textwidth}{!}{%
\begin{tikzpicture}[
    entity/.style={draw=darknavy, fill=accentblue!18, very thick, rectangle, rounded corners=4pt, align=left, inner sep=6pt, font=\scriptsize\ttfamily, minimum width=4cm},
    entityai/.style={entity, fill=accentpurple!18},
    entityfin/.style={entity, fill=accentgreen!18},
    entityaud/.style={entity, fill=accentorange!22},
    rel/.style={-, very thick, draw=darknavy!75},
    relai/.style={-, very thick, draw=accentpurple!85},
    relaud/.style={-, very thick, draw=accentorange!85}
]

% Эгнээ 1: Гол өгөгдөл (users, accounts, categories)
\node[entity]    (user) at (0,0)    {%
\textbf{users} \\
+ id PK\\
+ name, email UQ\\
+ password (bcrypt)\\
+ currency, push\_token};

\node[entityfin] (acc)  at (6,0)    {%
\textbf{accounts} \\
+ id PK\\
+ user\_id FK\\
+ name, type\\
+ balance, icon};

\node[entity]    (cat)  at (12,0)   {%
\textbf{categories} \\
+ id PK\\
+ name, icon, color\\
+ type};

% Эгнээ 2: Финанс модуль (budgets, transactions, scheduled_payments)
\node[entityfin] (bud)   at (0, -4) {%
\textbf{budgets} \\
+ id PK\\
+ user\_id, category\_id FK\\
+ amount, spent\\
+ month, year};

\node[entityfin] (tx)    at (6, -4) {%
\textbf{transactions} \\
+ id PK\\
+ user\_id FK\\
+ account\_id FK\\
+ category\_id FK\\
+ amount, date};

\node[entityfin] (sched) at (12, -4){%
\textbf{scheduled\_payments} \\
+ id PK\\
+ user\_id, account\_id FK\\
+ category\_id FK\\
+ frequency, next\_date};

% Эгнээ 3: AI модуль (ai_chats, ai_messages, ai_analyses)
\node[entityai] (chat)  at (0, -8)  {%
\textbf{ai\_chats} \\
+ id PK\\
+ user\_id FK\\
+ title};

\node[entityai] (msg)   at (6, -8)  {%
\textbf{ai\_messages} \\
+ id PK\\
+ chat\_id FK\\
+ role, content};

\node[entityai] (anal)  at (12, -8) {%
\textbf{ai\_analyses} \\
+ id PK\\
+ user\_id FK\\
+ filename, bank\_name\\
+ income, expenses\\
+ categories\_json};

% Эгнээ 4: Туслах (notifications, activity_logs)
\node[entityaud] (notif) at (0, -12) {%
\textbf{notifications} \\
+ id PK, user\_id FK\\
+ title, body, type\\
+ is\_read};

\node[entityaud] (log)   at (6, -12) {%
\textbf{activity\_logs} \\
+ id PK, user\_id FK\\
+ action, entity\\
+ details, status\\
+ ip\_address};

% --- Холболтууд (зөв өнцөгт шугам) ---

% Row 1: users → accounts; categories → accounts (горизонтальные)
\draw[rel] (user.east)  -- (acc.west);
\draw[rel] (cat.west)   -- (acc.east);

% users → budgets / tx (вниз)
\draw[rel] (user.south) -- (bud.north);
\draw[rel] (user.south east) |- ([yshift=4mm]tx.north);

% accounts → tx (вниз)
\draw[rel] (acc.south)  -- (tx.north);

% accounts → scheduled_payments (вниз с правой стороны)
\draw[rel] (acc.south east) |- ([yshift=4mm]sched.north);

% categories -> tx (доош, баруун талаас)
\draw[rel] (cat.south)  -- ([yshift=4mm]tx.north east);

% categories → budgets / scheduled (orthogonal)
\draw[rel] (cat.south west) |- ([yshift=4mm]bud.north east);
\draw[rel] (cat.south)  -- (sched.north);

% users → ai модуль (ai_chats зүүн талаар, ai_analyses баруун талаар)
\draw[relai] (user.west)  to[out=180, in=180] (chat.west);
\draw[relai] (user.east)  to[out=0,   in=0]   (anal.east);

% chat → msg
\draw[relai] (chat.east)  -- (msg.west);

% users → notifications / activity_logs
\draw[relaud] (user.west) to[out=180, in=180] (notif.west);
\draw[relaud] (user.south west) to[out=-90, in=180] (log.west);

\end{tikzpicture}
}
\caption{Entity-Relationship Diagram (ER диаграм) --- 11 хүснэгтийн foreign-key уялдаа}
\label{fig:erd}
\end{figure}

Хайрцагны өнгөний кодлол: цэнхэр --- цөм өгөгдөл, ногоон --- санхүүгийн модуль, нил ягаан --- AI модуль, улбар шар --- туслах (мэдэгдэл, аудит). Цэгцтэй байлгахын тулд цөөн тооны хэвтээ ба босоо шугам ашигласан.

\subsection{Foreign key хязгаарлалт}
Бүх FK-ууд ON DELETE CASCADE-тай тохируулагдсан. Хэрэв хэрэглэгч өөрийгөө устгах бол:
\begin{itemize}
    \item Бүх данс, гүйлгээ, төсөв, давтагдах төлбөр, чат, чат мессеж, шинжилгээ, мэдэгдэл, лог автоматаар устана.
    \item Гүйлгээ өөрөө устгагдвал зөвхөн данс л үлдэнэ (үлдэгдэл шинэчлэгдэнэ).
\end{itemize}

\section{Обьект хадгалалт (MinIO)}
\label{sec:minio}
Реляц өгөгдлийн сан (PostgreSQL) нь зөвхөн \emph{бүтэцлэгдсэн} мэдээллийг (тоо, текст, огноо, харилцаа) хадгална. Хэрэглэгчийн оруулсан \textbf{түүхий хоёртын файлууд} (PDF/XLSX/CSV банкны хуулга, профайл аватар) нь өөр зориулалттай хадгалалтын үйлчилгээ шаардана. Иймд систем \textbf{MinIO} \cite{minio} обьект хадгалалтыг сонгосон.

\subsection{Яагаад MinIO-г сонгосон бэ?}
\begin{itemize}
    \item \textbf{S3 API нийцтэй.} AWS S3-ийн PutObject, GetObject, ListObjects, presigned URL зэрэг бүх үндсэн дуудлагыг яг ижил аргаар дуудна. Ирээдүйд AWS S3, Backblaze B2, Cloudflare R2 рүү хөдөлгөөн хийхэд код өөрчлөгдөхгүй.
    \item \textbf{Нээлттэй эх + self-host.} AGPL лицензтэй, локал серверт суулгах боломжтой. Хувийн мэдээллийг өөрийн дэд бүтэцэд хадгалах боломжтой нь Монголын банкны хуулгын мэдрэмтгий байдалд тохиромжтой.
    \item \textbf{Өндөр гүйцэтгэл.} Go хэлээр бичигдсэн, нэг node 100,000+ обьект/секундийн уншилтыг даадаг.
    \item \textbf{Builtin шифрлэлт ба хувилбарлалт (versioning).} \texttt{SSE-S3} буюу \texttt{SSE-KMS}-ээр сервер тал шифрлэх боломжтой, мөн хувилбарлалт идэвхжүүлбэл санамсаргүй устгалаас сэргийлдэг.
    \item \textbf{Контейнерт уян хатан.} \texttt{docker-compose.yml}-д нэг шугамаар нэмнэ: \texttt{minio/minio:latest} дүрс, port 9000 (S3 API) + port 9001 (Web Console).
\end{itemize}

\subsection{Bucket-ийн зохион байгуулалт}
Систем хоёр bucket ашигладаг:

\begin{table}[ht]
\centering
\footnotesize
\renewcommand{\arraystretch}{1.25}
\rowcolors{2}{softgrey}{white}
\begin{tabular}{|p{3.5cm}|p{4cm}|p{5cm}|}
\hline
\rowcolor{darknavy}
\textcolor{white}{\textbf{Bucket нэр}} & \textcolor{white}{\textbf{Зориулалт}} & \textcolor{white}{\textbf{Файлын нэрийн загвар}} \\ \hline
\texttt{fintrack-statements}  & Банкны хуулга (PDF/XLSX/XLS/CSV)        & \texttt{statement\_<userID>\_<unix>.<ext>} \\ \hline
\texttt{fintrack-avatars}     & Профайл аватар зураг (PNG/JPG)          & \texttt{avatar\_<userID>\_<unix>.<ext>} \\ \hline
\end{tabular}
\caption{MinIO bucket-ийн зохион байгуулалт}
\label{tab:minio-buckets}
\end{table}

\subsection{Файл хадгалах урсгал}
\begin{enumerate}
    \item Хэрэглэгч \texttt{multipart/form-data}-аар сервер тал руу файл илгээнэ.
    \item Боловсруулагч (\texttt{ai\_analysis.go} эсвэл \texttt{auth.go:UploadAvatar}) файлыг түр зуурын буфферээс уншиж \texttt{minio-go} SDK-ийн \texttt{PutObject} аргаар bucket-д байршуулна.
    \item Үр дүнд буцсан \textbf{обьектын түлхүүр} (\texttt{key}, жишээ нь \texttt{statement\_42\_1717000000.pdf}) нь \texttt{ai\_analyses.filename} эсвэл \texttt{users.avatar} талбарт хадгалагдана. \emph{Бодит файл DB-д ордоггүй}, зөвхөн заагч.
    \item Хэрэглэгч талд харуулахын тулд \textbf{presigned URL} (15 мин TTL) үүсгэн дамжуулна. Энэ нь хэрэглэгч серверээр дамжихгүйгээр шууд MinIO-оос файлыг татах боломжийг олгож, серверийн ачааллыг бууруулдаг.
\end{enumerate}

\subsection{MinIO-ийн давуу тал ./uploads локал хавтсаас}
\begin{table}[ht]
\centering
\footnotesize
\renewcommand{\arraystretch}{1.3}
\rowcolors{2}{softgrey}{white}
\begin{tabular}{|p{3.5cm}|p{4.2cm}|p{4.5cm}|}
\hline
\rowcolor{darknavy}
\textcolor{white}{\textbf{Шинж чанар}} & \textcolor{white}{\textbf{Локал \texttt{./uploads}}} & \textcolor{white}{\textbf{MinIO bucket}} \\ \hline
Хэвтээ масштабжилт & Боломжгүй (сервер бүрд тус тусдаа файл) & Олон backend нэгж нэг bucket-аас уншина \\ \hline
Шифрлэлт            & Гар тохируулга                          & SSE-S3 builtin \\ \hline
Хувилбарлалт         & Байхгүй                                 & Versioning идэвхжүүлбэл санамсаргүй устгал сэргээгдэнэ \\ \hline
Гадаад сервис рүү шилжих & Кодыг бичих ёстой                      & S3-нийцтэй тул AWS S3, B2, R2 рүү яг ижил код-оор \\ \hline
Backup, replication  & Локал диск-тэй уях                      & Хэд хэдэн node-д хуулбар (erasure code) \\ \hline
Хандах хяналт         & OS файлын зөвшөөрөл                     & IAM-style bucket policy, presigned URL \\ \hline
\end{tabular}
\caption{Локал файл систем ба MinIO-ийн харьцуулалт}
\label{tab:fs-vs-minio}
\end{table}

\section{REST API цэгийн зохиомж}
Систем нийт \textbf{45 API цэгтэй} бөгөөд RESTful зарчмын \cite{rest} дагуу боловсруулсан:
\begin{itemize}
    \item Нөөц төвтэй URL (\texttt{/transactions}, \texttt{/budgets}),
    \item HTTP аргуудыг семантик зориулалтаар нь ашигласан (GET --- унших, POST --- үүсгэх, PUT --- шинэчлэх, DELETE --- устгах),
    \item Хариулт нь стандарт JSON формат, HTTP төлвийн код зөв ашигласан (200, 201, 400, 401, 404, 409, 500),
    \item Хуудаслалт нь \texttt{?page=1\&limit=20} асуулгын параметрээр.
\end{itemize}

\begin{table}[ht]
\centering
\footnotesize
\renewcommand{\arraystretch}{1.2}
\rowcolors{2}{softgrey}{white}
\begin{tabular}{|c|p{4.2cm}|p{6cm}|}
\hline
\rowcolor{darknavy}
\textcolor{white}{\textbf{Method}} & \textcolor{white}{\textbf{Path}} & \textcolor{white}{\textbf{Зориулалт}} \\ \hline
POST   & \texttt{/auth/register}            & Бүртгэл \\ \hline
POST   & \texttt{/auth/login}               & Имэйл + нууц үгээр нэвтрэлт \\ \hline
POST   & \texttt{/auth/forgot-password}     & Нууц үг сэргээх хүсэлт \\ \hline
GET    & \texttt{/auth/social/providers}    & Нийгмийн нэвтрэлтийн нийлүүлэгчийн жагсаалт \\ \hline
POST   & \texttt{/auth/social/login}        & Google/Apple/Facebook нэвтрэлт \\ \hline
GET    & \texttt{/profile}                  & Профайл харах \\ \hline
PUT    & \texttt{/profile}                  & Профайл шинэчлэх \\ \hline
PUT    & \texttt{/profile/password}         & Нууц үг солих \\ \hline
POST   & \texttt{/profile/avatar}           & Аватар зураг солих \\ \hline
GET    & \texttt{/dashboard}                & Хяналтын самбарын үзүүлэлт \\ \hline
GET    & \texttt{/expenses/summary}         & Зарлагын нэгтгэл \\ \hline
GET    & \texttt{/statistics}               & Хугацаа дамнасан статистик \\ \hline
POST/GET/DELETE & \texttt{/transactions}    & Гүйлгээний CRUD (4) \\ \hline
POST/GET/PUT/DELETE & \texttt{/budgets}     & Сарын төсвийн CRUD (4) \\ \hline
POST/GET/PUT/DELETE & \texttt{/accounts}    & Дансны CRUD (4) \\ \hline
GET, POST & \texttt{/categories}            & Ангилал унших, нэмэх (2) \\ \hline
POST/GET/DELETE/PUT & \texttt{/scheduled-payments} & Давтагдах төлбөрийн CRUD (4) \\ \hline
PUT    & \texttt{/scheduled-payments/:id/toggle} & Идэвхтэй/идэвхгүй болгох \\ \hline
POST   & \texttt{/ai/chat}                  & AI чат руу мессеж илгээх \\ \hline
POST/GET/DELETE & \texttt{/ai/chats}        & Чат сесс-ийн CRUD (4) \\ \hline
\textbf{POST}   & \textbf{\texttt{/ai/analysis}} & \textbf{Банкны хуулга AI шинжилгээ} \\ \hline
GET, DELETE & \texttt{/ai/analyses}         & Шинжилгээний түүх (3) \\ \hline
POST   & \texttt{/import/statement}         & Текст-only хуулга import \\ \hline
POST   & \texttt{/push-token}               & Expo push токен бүртгэл \\ \hline
GET, PUT  & \texttt{/notifications}         & Мэдэгдэл удирдах (4) \\ \hline
GET    & \texttt{/activity-logs}            & Аудит лог жагсаалт \\ \hline
GET    & \texttt{/activity-logs/summary}    & Аудит нэгтгэлийн статистик \\ \hline
\end{tabular}
\caption{REST API-гийн гол цэгүүдийн жагсаалт}
\label{tab:endpoints}
\end{table}

\section{AI шинжилгээний хоолой}
Системийн хамгийн нарийн хэсэг нь банкны хуулгыг автомат байдлаар бүтэцтэй форматруу хөрвүүлж AI-д тэжээх хоолой юм.

\begin{figure}[ht]
\centering
\begin{tikzpicture}[
    pipestep/.style={rectangle, rounded corners, draw=darknavy, very thick, minimum width=2.4cm, minimum height=1cm, align=center, font=\scriptsize},
    s1/.style={pipestep, fill=accentblue!25},
    s2/.style={pipestep, fill=accentteal!25},
    s3/.style={pipestep, fill=accentgreen!25},
    s4/.style={pipestep, fill=accentyellow!40},
    s5/.style={pipestep, fill=accentpurple!25},
    s6/.style={pipestep, fill=accentorange!25},
    arrow/.style={-{Stealth[length=3mm]}, very thick, draw=darknavy}
]

\node[s1] (up)   at (0,0)      {1. Байршуулах\\PDF/CSV/XLSX};
\node[s2] (sv)   at (3.5,0)    {2. Файл\\хадгалах};
\node[s3] (ps)   at (7,0)      {3. Задлан\\унших};
\node[s4] (cb)   at (10.5,0)   {4. Ангиллын\\задаргаа};
\node[s5] (ai)   at (10.5,-3)  {5. AI\\хураангуй};
\node[s6] (db)   at (3.5,-3)   {6. DB-д\\хадгалах};
\node[s1] (out)  at (0,-3)     {7. JSON\\хариулт};

\draw[arrow] (up.east) -- (sv.west);
\draw[arrow] (sv.east) -- (ps.west);
\draw[arrow] (ps.east) -- (cb.west);
\draw[arrow] (cb.south) -- (ai.north);
\draw[arrow] (ai.west) -- (db.east);
\draw[arrow] (db.west) -- (out.east);

\end{tikzpicture}
\caption{Хуулга → AI шинжилгээ урсгал (7 алхам)}
\label{fig:ai-урсгал}
\end{figure}

\subsection{Алхам бүрийн дэлгэрэнгүй}
\begin{enumerate}
    \item \textbf{Файл байршуулах.} Хэрэглэгч \texttt{multipart/form-data} ашиглан POST \texttt{/api/v1/ai/analysis} руу файл илгээнэ. PDF, XLSX, XLS, CSV-г зөвшөөрнө.
    \item \textbf{Файл хадгалах (MinIO).} Сервер файлыг MinIO обьект хадгалалтын \texttt{fintrack-statements} bucket-д \texttt{statement\_<userID>\_<unix>.<ext>} нэрээр байршуулна. Бүх raw файл (PDF/XLSX/CSV банкны хуулга болон профайл аватар зураг) MinIO дотор хадгалагдаж, гүйцэтгэл, тогтвортой байдал, хувилбарлалт (versioning) болон уян хатан масштабжилтыг хангадаг. Үйлчилгээ нь S3 API-тай нийцтэй тул ирээдүйд AWS S3, Backblaze B2 зэрэг бусад нийлүүлэгч рүү амархан шилжих боломжтой.
    \item \textbf{Задлан унших.} Python хуулга задлагч ажиллаж байвал HTTP-ээр шууд дамжуулагдана; үгүй бол Go-ийн нөөц задлагч ашиглагдана. Хариулт нь \texttt{ParsedStatement} бүтэц юм.
    \item \textbf{Ангиллын задаргаа.} Хэрэв задлагчаас задаргаа ирээгүй бол гүйлгээний жагсаалтаас ангиллаар бүлэглэх функц ажиллаж хувь, тоог тооцоолно.
    \item \textbf{AI хураангуй.} \texttt{ai\_provider.go} модулийн \texttt{generateAIResponse} функц нь үндсэн нийлүүлэгч (анхдагч \texttt{gemini-2.5-flash}) рүү системийн зааварчилгаа болон JSON контекстийг илгээж, амжилтгүй бол OpenRouter руу автомат нөөц шилжилт хийж хүн уншигдахуйц хураангуй болон 3-5 зөвлөмжийн жагсаалт буцаана.
    \item \textbf{Өгөгдлийн санд хадгалах.} Бүх үр дүнг \texttt{ai\_analyses} хүснэгтэд хадгална (categories\_json, transactions\_json, recommendations\_json).
    \item \textbf{Хариулт.} Хэрэглэгч тал руу \texttt{AIAnalysisResponse} JSON-оор бүх дугаар, графикийн өгөгдөл, зөвлөмж, AI хураангуй буцаана.
\end{enumerate}

\subsection{AI чатын алхамчилсан гүйцэтгэл}

\begin{figure}[ht]
\centering
\begin{tikzpicture}[
    actor/.style={draw=darknavy, fill=accentteal!30, thick, rectangle, rounded corners, minimum width=2cm, minimum height=0.7cm, font=\footnotesize},
    msg/.style={->, thick, draw=darknavy, font=\scriptsize},
    ret/.style={->, thick, draw=accentgreen, dashed, font=\scriptsize}
]

\node[actor] (u)  at (0, 0)  {Хэрэглэгч};
\node[actor] (fe) at (3, 0)  {Хэрэглэгч тал};
\node[actor] (be) at (6, 0)  {Сервер тал};
\node[actor] (db) at (9, 0)  {DB};
\node[actor] (ai) at (12, 0) {AI / LLM};

\foreach \x in {0,3,6,9,12} {
    \draw[dashed, thick, gray] (\x, -0.4) -- (\x, -7);
}

\draw[msg] (0,-1) -- node[above]{асуулт бичих} (3,-1);
\draw[msg] (3,-1.7) -- node[above]{POST /ai/chat} (6,-1.7);
\draw[msg] (6,-2.4) -- node[above]{контекст (данс/сар) Tx} (9,-2.4);
\draw[ret] (9,-3.0) -- node[above]{Tx + balance} (6,-3.0);
\draw[msg] (6,-3.7) -- node[above]{зааварчилгаа + контекст} (12,-3.7);
\draw[ret] (12,-4.4) -- node[above]{LLM хариулт} (6,-4.4);
\draw[msg] (6,-5.1) -- node[above]{AIMessage хадгалах} (9,-5.1);
\draw[ret] (6,-5.8) -- node[above]{200 OK} (3,-5.8);
\draw[ret] (3,-6.5) -- node[above]{дүрслэх} (0,-6.5);

\end{tikzpicture}
\caption{Sequence Diagram (Дарааллын диаграм) --- AI чат}
\label{fig:seq-chat}
\end{figure}

\subsection{Гүйлгээ нэмэх үйл ажиллагааны диаграм}
Дараах үйл ажиллагааны диаграм нь хэрэглэгч шинэ гүйлгээ нэмэхэд системд хийгдэх алхмууд, шийдвэрийн салаалт, хэвийн ба алдаатай гарцыг тодорхой үзүүлж байна.

\begin{figure}[ht]
\centering
\begin{tikzpicture}[
    >={Stealth[length=3mm]},
    start/.style={circle, draw=darknavy, fill=darknavy, minimum size=6mm, inner sep=0pt},
    finish/.style={circle, draw=darknavy, very thick, minimum size=7mm, inner sep=0pt, double, double distance=1.2pt},
    enderr/.style={circle, draw=accentorange!85!black, fill=accentorange!85!black, very thick, minimum size=4mm, inner sep=0pt, path picture={\draw[white, line width=0.6pt] (path picture bounding box.north west) -- (path picture bounding box.south east); \draw[white, line width=0.6pt] (path picture bounding box.north east) -- (path picture bounding box.south west);}},
    act/.style={rectangle, rounded corners=4pt, draw=darknavy, very thick, fill=accentblue!22, align=center, font=\scriptsize, minimum width=4.5cm, minimum height=0.9cm},
    actfin/.style={act, fill=accentgreen!22},
    actai/.style={act, fill=accentpurple!22},
    decision/.style={diamond, draw=darknavy, very thick, fill=accentyellow!45, align=center, font=\scriptsize, aspect=2, inner sep=1pt},
    error/.style={rectangle, rounded corners=4pt, draw=accentorange!80!black, very thick, fill=accentorange!35, align=center, font=\scriptsize, minimum width=3.5cm, minimum height=0.9cm},
    flow/.style={->, very thick, draw=darknavy},
    flowerr/.style={->, very thick, draw=accentorange!85!black, dashed},
    branchlbl/.style={font=\scriptsize\bfseries, inner sep=2pt}
]

\node[start] (s) at (0,0) {};
\node[act]   (form) at (0,-1.2) {Хэрэглэгч гүйлгээний\\ маягт бөглөх};
\node[act]   (submit) at (0,-2.7) {POST /api/v1/transactions\\ илгээх};

\node[decision] (jwt)   at (0,-4.4)  {JWT хүчинтэй\\ юу?};
\node[error]    (e1)    at (5.5,-4.4){401 Unauthorized};
\node[enderr]   (ef1)   at (9.5,-4.4){};

\node[decision] (valid) at (0,-6.4)  {Талбарууд\\ зөв үү?};
\node[error]    (e2)    at (5.5,-6.4){400 Bad Request};
\node[enderr]   (ef2)   at (9.5,-6.4){};

\node[actfin]   (lookup) at (0,-8.2) {Account, Category\\ DB-ээс олох};

\node[decision] (own)   at (0,-10.0) {Хэрэглэгчид\\ хамаарах уу?};
\node[error]    (e3)    at (5.5,-10.0){403 Forbidden};
\node[enderr]   (ef3)   at (9.5,-10.0){};

\node[actfin]   (insert) at (0,-11.8) {transactions хүснэгтэд\\ оруулах (INSERT)};
\node[actfin]   (bal)    at (0,-13.3) {Дансны үлдэгдэл\\ дахин тооцох};
\node[actfin]   (budget) at (0,-14.8) {Хамаарагч төсвийн\\ Spent шинэчлэх};
\node[actai]    (audit)  at (0,-16.3) {activity\_logs-д\\ үйлдэл бүртгэх};
\node[act]      (resp)   at (0,-17.8) {201 Created + JSON\\ хариулт буцаах};
\node[finish]   (f)      at (0,-19.2) {};

% Үндсэн (амжилтын) урсгал
\draw[flow] (s) -- (form);
\draw[flow] (form) -- (submit);
\draw[flow] (submit) -- (jwt);
\draw[flow] (jwt.south)   -- node[branchlbl, right]{Тийм} (valid.north);
\draw[flow] (valid.south) -- node[branchlbl, right]{Тийм} (lookup.north);
\draw[flow] (lookup) -- (own);
\draw[flow] (own.south)   -- node[branchlbl, right]{Тийм} (insert.north);
\draw[flow] (insert) -- (bal);
\draw[flow] (bal)    -- (budget);
\draw[flow] (budget) -- (audit);
\draw[flow] (audit)  -- (resp);
\draw[flow] (resp)   -- (f);

% Алдааны салаа гарцууд (тус бүрдээ өөрийн төгсгөлтэй)
\draw[flowerr] (jwt.east)   -- node[branchlbl, above]{Үгүй} (e1.west);
\draw[flowerr] (e1.east)    -- (ef1);
\draw[flowerr] (valid.east) -- node[branchlbl, above]{Үгүй} (e2.west);
\draw[flowerr] (e2.east)    -- (ef2);
\draw[flowerr] (own.east)   -- node[branchlbl, above]{Үгүй} (e3.west);
\draw[flowerr] (e3.east)    -- (ef3);

\end{tikzpicture}
\caption{Activity Diagram (Үйл ажиллагааны диаграм) --- гүйлгээ нэмэх}
\label{fig:activity-tx}
\end{figure}

\textbf{Тайлбар.} Шар ромбоор шийдвэрийн салаалтыг (JWT баталгаажуулалт, оролтын баталгаажуулалт, эзэмшил шалгалт), цэнхэр-ногоон тэгш өнцөгтөөр төлөв шинэчлэх үйлдлүүдийг, улбар шар тэгш өнцөгтөөр алдааны гарцыг харуулав. Алдаа бүр өөрийн ($\otimes$) төгсгөлийн тэмдэгтэй --- гүйцэтгэл тухайн алхамд зогсож, хэрэглэгч рүү HTTP төлөв буцаагдана. Хэвийн урсгалд 7 алхам дамжин, эцэст нь 201 Created хариу буцна.

\section{Хэрэглэгч талын зохиомж}

\subsection{Дэлгэцийн жагсаалт}
React Native гар утасны программ нь нийт \textbf{17 дэлгэцтэй} ба \texttt{src/screens/} хавтсанд төвлөрсөн.

\begin{table}[ht]
\centering
\footnotesize
\renewcommand{\arraystretch}{1.2}
\rowcolors{2}{softgrey}{white}
\begin{tabular}{|p{4.5cm}|p{8cm}|}
\hline
\rowcolor{darknavy}
\textcolor{white}{\textbf{Дэлгэц}} & \textcolor{white}{\textbf{Зориулалт}} \\ \hline
OnboardingScreen      & Анхны нээлтийн угтах дэлгэц. \\ \hline
LoginScreen           & Нэвтрэх. \\ \hline
RegisterScreen        & Бүртгүүлэх. \\ \hline
HomeScreen            & Хяналтын самбар, үлдэгдэл, гүйлгээ. \\ \hline
TransactionsScreen    & Бүх гүйлгээний жагсаалт, шүүлт. \\ \hline
AddTransactionScreen  & Шинэ гүйлгээ. \\ \hline
ExpensesScreen        & Ангилал бүрийн зарлага + тойрог график. \\ \hline
StatisticsScreen      & Хугацаагаар (өдөр/долоо/сар/жил) багана график. \\ \hline
BudgetScreen          & Сарын төсөв + прогресс. \\ \hline
AccountsScreen        & Дансны жагсаалт, CRUD. \\ \hline
ScheduledPaymentsScreen & Давтагдах төлбөр. \\ \hline
AIChatScreen          & AI санхүүгийн зөвлөгөөч чат. \\ \hline
DataImportScreen      & Банкны хуулга байршуулах, AI шинжилгээ. \\ \hline
NotificationsScreen   & Push мэдэгдлүүд. \\ \hline
ProfileScreen         & Профайл засах. \\ \hline
SettingsScreen        & Систем тохиргоо. \\ \hline
PrivacyScreen         & Нууцлал, политик. \\ \hline
\end{tabular}
\caption{Хэрэглэгч талын дэлгэцийн жагсаалт}
\label{tab:screens}
\end{table}

\subsection{Хэрэглэгч талын кодын зохион байгуулалт}

\begin{table}[ht]
\centering
\footnotesize
\renewcommand{\arraystretch}{1.2}
\rowcolors{2}{softgrey}{white}
\begin{tabular}{|l|p{8.5cm}|}
\hline
\rowcolor{darknavy}
\textcolor{white}{\textbf{Хавтас}} & \textcolor{white}{\textbf{Зориулалт}} \\ \hline
\texttt{src/screens/}      & 17 дэлгэц. \\ \hline
\texttt{src/navigation/}   & Stack + Bottom Tab navigator. \\ \hline
\texttt{src/services/}     & Axios HTTP үйлчилгээ, төрөлтэй. \\ \hline
\texttt{src/context/}      & AuthContext (нэвтрэх/гарах/бүртгүүлэх). \\ \hline
\texttt{src/types/}        & TypeScript интерфейс, type. \\ \hline
\texttt{src/config/}       & Сүлжээний байршуулалт (төвлөрсөн тохиргоо). \\ \hline
\texttt{tailwind.config.js} & Харанхуй загварын утгууд. \\ \hline
\texttt{global.css}        & NativeWind base. \\ \hline
\end{tabular}
\caption{Хэрэглэгч талын кодын зохион байгуулалт}
\label{tab:fe-layout}
\end{table}

\subsection{Дизайн зарчим}
Дэлгэцийн дизайн нь дараах зарчмыг баримталсан:
\begin{itemize}
    \item \textbf{Харанхуй загвар.} Хэрэглэгчийн нүдийг бага ачаална, орчин үеийн гар утасны программуудын нийтлэг стандарт.
    \item \textbf{Гар утас төвтэй (mobile-first) хандлага.} Гар утасны дэлгэцэнд зориулсан, илүүдэл том элементгүй.
    \item \textbf{Картын систем.} Мэдээллийг бие даасан карт (card)-ууд хэлбэрээр харуулна.
    \item \textbf{Тогтмол өнгөний код.} Орлого --- ногоон, зарлага --- улаан, AI --- ягаан.
    \item \textbf{Микро харилцан үйлчлэл.} Товчлуур дарах, гүйлгээ нэмэх үед анимац үзүүлнэ.
\end{itemize}

\subsection{Гол дэлгэцийн зохиомж}
Доорх диаграмаас гол дэлгэцүүдийн схем зохион байгуулалтыг харж болно.

\begin{figure}[ht]
\centering
\begin{tikzpicture}[
    phone/.style={draw=darknavy, very thick, rounded corners=8pt, fill=darknavy, minimum width=4cm, minimum height=7cm},
    screen/.style={draw=midgrey, fill=white, minimum width=3.6cm, minimum height=6.4cm, font=\scriptsize, align=center},
    head/.style={fill=accentblue, draw=accentblue, text=white, font=\scriptsize\bfseries, minimum width=3.6cm, minimum height=0.7cm, align=center},
    section/.style={fill=softgrey, draw=midgrey, font=\scriptsize, minimum width=3.4cm, minimum height=0.9cm, align=center}
]

% Phone 1: Home
\node[phone] (p1) at (0,0) {};
\node[head] at (0, 3.0) {Нүүр / Самбар};
\node[section, fill=accentpurple!30] at (0, 1.9) {Total Balance \\ \textbf{MNT 4.25M}};
\node[section, fill=accentgreen!30] at (-0.8, 0.7) {Income};
\node[section, fill=accentred!30] at (0.8, 0.7) {Expense};
\node[section] at (0, -0.5) {Recent Tx};
\node[section] at (0, -1.6) {Quick add};
\node[section, fill=accentyellow!40] at (0, -2.7) {AI Insight};

% Phone 2: AI Chat
\node[phone] (p2) at (5,0) {};
\node[head, fill=accentpurple] at (5, 3.0) {AI Advisor};
\node[section, fill=accentpurple!15, minimum height=0.6cm] at (5, 2.1) {Сар бүр зарлага \\ яагаад өсдөг вэ?};
\node[section, fill=accentteal!20, minimum height=1.2cm] at (5, 1.0) {AI: Энэ сард \\ хоол, тээвэр ангилалд...};
\node[section, fill=accentpurple!15, minimum height=0.6cm] at (5, -0.2) {Хадгаламж яаж \\ ихэсгэх вэ?};
\node[section, fill=accentteal!20, minimum height=1.2cm] at (5, -1.3) {AI: Subscription-уудаа \\ шалгахыг санал болгоё...};
\node[section, fill=midgrey] at (5, -2.7) {[ Type a message ]};

% Phone 3: AI analysis
\node[phone] (p3) at (10,0) {};
\node[head, fill=accentorange] at (10, 3.0) {Statement Analysis};
\node[section, fill=accentblue!20] at (10, 2.0) {Period: 2026.04};
\node[section, fill=accentgreen!20, minimum height=0.7cm] at (10, 1.0) {Income: MNT 1.8M};
\node[section, fill=accentred!20, minimum height=0.7cm] at (10, 0.0) {Expense: MNT 1.4M};
\node[section, fill=accentyellow!30, minimum height=0.8cm] at (10, -1.1) {Ангилал \\ тойрог график};
\node[section, fill=accentpurple!20, minimum height=1cm] at (10, -2.4) {AI зөвлөмж: \\ 3 tip};

\end{tikzpicture}
\caption{Үндсэн гурван дэлгэцийн зохиомж}
\label{fig:mockup}
\end{figure}

\section{Аюулгүй байдал ба өгөгдлийн хамгаалалт}

\subsection{Баталгаажуулалт ба эрх олголт}
\begin{itemize}
    \item Бүх \texttt{/api/v1/*} API цэг (\texttt{auth/*}-аас бусад) JWT шаарддаг.
    \item Токен нь HMAC-SHA256-аар гарын үсэгтэй (\texttt{JWT\_SECRET} орчны хувьсагч).
    \item TTL = 72 цаг (3 хоног).
    \item Дундын давхарга нь \texttt{Authorization} толгой хэсгийг шалгаж user\_id-г Gin context-д тавина.
    \item Токен хүчингүй болсон бол клиент сэргээх (refresh) механизмгүй, дахин нэвтрэх ёстой.
\end{itemize}

\subsection{Нууц үгийн аюулгүй байдал}
\begin{itemize}
    \item Bcrypt cost 10 (анхдагч) ашиглана.
    \item Ил текст нууц үг ямар ч лог, хариуд буцдаггүй (\texttt{json:\textquotedbl-\textquotedbl}).
    \item Хамгийн багадаа 6 тэмдэгт шаардлага.
\end{itemize}

\subsection{SQL injection хамгаалалт}
GORM нь параметрлэсэн асуулгыг анхдагчаар ашигладаг тул SQL injection халдлага хориглогдсон. Бүх боловсруулагчид \texttt{Where(...)}-ийг хоёр параметртэй хэлбэрээр, утгын байрлуулагч ашиглан дуудсан.

\subsection{CORS}
Хэрэглэгч тал сервер тал руу хандах боломжийг \texttt{middleware/cors.go}-аас зөвшөөрсөн. Үйлдвэрлэлийн орчинд тодорхой origin-ыг зөвшөөрөгдсөн жагсаалтад оруулна.

\subsection{HTTPS}
Үйлдвэрлэлийн орчинд reverse proxy (nginx эсвэл Caddy) дээр TLS termination хийнэ. Let's Encrypt-ийн үнэгүй сертификат ашигладаг.

\section{Контейнержуулалт ба ажиллах орчин}
\textbf{docker-compose.yml} нь 4 үйлчилгээг зэрэг ажиллуулна: \texttt{postgres} (port 5433 → 5432), \texttt{minio} (port 9000 API, 9001 console), \texttt{parser} (Python FastAPI хуулга задлагч, port 8000) ба \texttt{backend} (port 8080). Сервер талын Dockerfile нь олон шатлалт бүтээлттэй:
\begin{enumerate}
    \item \textbf{Бүтээгч шат.} \texttt{golang:1.25-alpine}-ээр Go binary-г \texttt{CGO\_ENABLED=0} тохиргоотой хөрвүүлнэ.
    \item \textbf{Эцсийн шат.} \texttt{alpine:3.19} дээр зөвхөн binary + ca-certificates + tzdata байрлана. Эцсийн дүрс \textasciitilde{}20МБ хэмжээтэй.
\end{enumerate}

\begin{table}[ht]
\centering
\footnotesize
\renewcommand{\arraystretch}{1.25}
\rowcolors{2}{softgrey}{white}
\begin{tabular}{|l|p{4cm}|p{5cm}|}
\hline
\rowcolor{darknavy}
\textcolor{white}{\textbf{Үйлчилгээ}} & \textcolor{white}{\textbf{Дүрс}} & \textcolor{white}{\textbf{Тэмдэглэл}} \\ \hline
\texttt{postgres} & \texttt{postgres:16-alpine}            & port 5432→5433, healthcheck. \\ \hline
\texttt{minio}    & \texttt{minio/minio:latest}            & port 9000 (S3 API), 9001 (console), volume \texttt{minio\_data}. \\ \hline
\texttt{parser}   & \texttt{./parser\_service}-ээс бүтээх  & port 8000, FastAPI хуулга задлагч, healthcheck. \\ \hline
\texttt{backend}  & \texttt{./backend}-ээс бүтээх          & port 8080, \texttt{depends\_on}: postgres, minio, parser. \\ \hline
\end{tabular}
\caption{Docker Compose үйлчилгээний тохиргоо}
\label{tab:docker}
\end{table}

\section{Хэрэгжүүлэлтийн жишээ кодууд}

\subsection{Хэрэглэгч бүртгэлийн боловсруулагч}
\begin{lstlisting}[language=Go, basicstyle=\ttfamily\scriptsize, caption={Бүртгэлийн боловсруулагч: bcrypt, JWT, анхдагч данс}]
func (h *AuthHandler) Register(c *gin.Context) {
    var req models.RegisterRequest
    if err := c.ShouldBindJSON(&req); err != nil {
        c.JSON(400, gin.H{"error": err.Error()}); return
    }
    var existing models.User
    if err := h.DB.Where("email = ?", req.Email).First(&existing).Error; err == nil {
        c.JSON(409, gin.H{"error": "Email already registered"}); return
    }
    hash, _ := bcrypt.GenerateFromPassword([]byte(req.Password), bcrypt.DefaultCost)
    user := models.User{Name: req.Name, Email: req.Email, Password: string(hash)}
    h.DB.Create(&user)
    h.DB.Create(&models.Account{UserID: user.ID, Name: "Main Account", Type: "bank"})
    token, _ := h.generateToken(user.ID)
    c.JSON(201, models.AuthResponse{Token: token, User: user})
}
\end{lstlisting}

\section{Олон нийлүүлэгчтэй AI архитектур}
Дипломын ажлын хамгийн өвөрмөц техникийн шийдэл бол \textbf{олон нийлүүлэгчтэй AI}-ийн архитектур юм. Энэ нь нэг талаас Монгол Улсаас Gemini API руу шууд хандах боломжгүй газарзүйн хязгаарлалтыг шийддэг, нөгөө талаас аль нэг үйлчилгээ доголдсон тохиолдолд системийг зогсоохгүй байх баталгааг хангадаг.

\subsection{Нийлүүлэгчийн дараалал}
\begin{enumerate}
    \item \textbf{Үндсэн нийлүүлэгч:} Google Gemini --- хэрэв \texttt{AI\_API\_KEY} тохируулагдсан бол эхлээд оролдоно.
    \item \textbf{Нөөц нийлүүлэгч:} OpenRouter --- \texttt{OPENROUTER\_API\_KEY} тохируулагдсан үед автоматаар идэвхжинэ.
\end{enumerate}

\subsection{Нөөц шилжилтийн нөхцлүүд}
Нийлүүлэгчийг солих шийдвэр дараах алдаанд тулгуурлана:
\begin{itemize}
    \item HTTP 429 (rate limit) --- үнэгүй түвшний өдрийн квот дууссан.
    \item HTTP 5xx --- үйлчилгээний серверийн алдаа.
    \item HTTP 451 эсвэл "User location is not supported" текст --- газарзүйн хязгаарлалт.
    \item Хариу хоосон ирэх.
\end{itemize}

\subsection{OpenRouter-ийн үнэгүй загварын жагсаалтыг автоматаар шинэчлэх}
OpenRouter-ийн үнэгүй загварын нэр тогтмол солигддог тул кодын статик жагсаалт хадгалах нь хурдан хуучирдаг. Шийдэл:
\begin{itemize}
    \item Систем 1 цаг тутамд \texttt{GET https://openrouter.ai/api/v1/models}-аас бүх загварын жагсаалтыг татна.
    \item \texttt{pricing.prompt = 0} болон \texttt{pricing.completion = 0} тохиргоотой загварыг үнэгүй гэж тэмдэглэнэ.
    \item Кэшийг \texttt{sync.Mutex}-ээр хамгаалсан санах ой дахь map-д хадгална.
    \item Тохиргооны \texttt{OPENROUTER\_MODEL} жагсаалт амжилтгүй болсон тохиолдолд автоматаар татсан жагсаалт руу шилжинэ.
\end{itemize}

\begin{table}[ht]
\centering
\footnotesize
\renewcommand{\arraystretch}{1.25}
\rowcolors{2}{softgrey}{white}
\begin{tabular}{|p{6.5cm}|p{2.5cm}|p{2.6cm}|}
\hline
\rowcolor{darknavy}
\textcolor{white}{\textbf{Загвар}} & \textcolor{white}{\textbf{Үйлдвэрлэгч}} & \textcolor{white}{\textbf{Параметр}} \\ \hline
\texttt{meta-llama/llama-3.3-70b-instruct:free}    & Meta       & 70B \\ \hline
\texttt{qwen/qwen-2.5-72b-instruct:free}           & Alibaba    & 72B \\ \hline
\texttt{google/gemma-2-9b-it:free}                 & Google     & 9B  \\ \hline
\texttt{mistralai/mistral-7b-instruct:free}        & Mistral AI & 7B  \\ \hline
\texttt{meta-llama/llama-3.2-3b-instruct:free}     & Meta       & 3B  \\ \hline
\texttt{deepseek/deepseek-chat:free}               & DeepSeek   & 67B \\ \hline
\end{tabular}
\caption{Анхдагч OpenRouter үнэгүй загварын жагсаалт}
\label{tab:or-models}
\end{table}

\subsection{AI / LLM үйлчилгээ дуудах фрагмент}
\begin{lstlisting}[language=Go, basicstyle=\ttfamily\scriptsize, caption={Контекстэд тулгуурласан LLM хариу (gemini-2.5-flash)}]
client, _ := genai.NewClient(context.Background(), option.WithAPIKey(apiKey))
defer client.Close()
model := client.GenerativeModel("gemini-2.5-flash")
prompt := fmt.Sprintf("You are a personal finance advisor. Context: %s. Question: %s",
    contextJSON, userMessage)
resp, _ := model.GenerateContent(ctx, genai.Text(prompt))
return resp.Candidates[0].Content.Parts[0].(genai.Text)
\end{lstlisting}

\subsection{React Native AuthContext}
\begin{lstlisting}[language=,basicstyle=\ttfamily\scriptsize, caption={AuthContext нэвтрэх функц (TypeScript)}]
const login = async (email: string, password: string) => {
  const { data } = await api.post('/auth/login', { email, password });
  await AsyncStorage.setItem('token', data.token);
  setUser(data.user);
};
\end{lstlisting}

\section{Хэрэгжүүлэлтийн тоон үзүүлэлт}

\begin{table}[ht]
\centering
\footnotesize
\renewcommand{\arraystretch}{1.25}
\rowcolors{2}{softgrey}{white}
\begin{tabular}{|p{7cm}|c|}
\hline
\rowcolor{darknavy}
\textcolor{white}{\textbf{Үзүүлэлт}} & \textcolor{white}{\textbf{Утга}} \\ \hline
Сервер талын Go код (мөр)         & $\approx$4{,}900 \\ \hline
Хэрэглэгч талын TS/TSX (мөр)      & $\approx$6{,}000 \\ \hline
Боловсруулагч файл                & 12 \\ \hline
REST API цэг                      & 45 \\ \hline
Өгөгдлийн сангийн хүснэгт         & 11 \\ \hline
React Native дэлгэц               & 17 \\ \hline
Эх ангилал (JSON)                 & 17 (12 зарлага + 5 орлого) \\ \hline
Docker контейнер                  & 4 (postgres, minio, parser, backend) \\ \hline
MinIO bucket                       & 2 (\texttt{fintrack-statements}, \texttt{fintrack-avatars}) \\ \hline
Дэмждэг файлын формат             & 3 (PDF, XLSX, CSV) \\ \hline
Дэмжих банкны формат              & 5 (Хаан, Голомт, ХХБ, Хас, Төрийн) \\ \hline
AI үндсэн загвар                  & Gemini 2.5 Flash \\ \hline
AI нөөц нийлүүлэгч                  & OpenRouter (6+ үнэгүй загвар) \\ \hline
\end{tabular}
\caption{Хэрэгжүүлсэн системийн тоон үзүүлэлтүүд}
\label{tab:metrics}
\end{table}

\textbf{Дүгнэлт.} Энэ бүлэгт системийн архитектур, сервер талын кодын зохион байгуулалт, загваруудын class диаграм, өгөгдлийн сангийн ER загвар, MinIO обьект хадгалалтын дизайн, REST API цэгүүд, AI шинжилгээний долоон алхамт урсгал, дарааллын диаграм, гүйлгээ нэмэх үйл ажиллагааны диаграм, хэрэглэгч талын дэлгэц, аюулгүй байдлын механизм болон хэрэгжүүлэлтийн тоон үзүүлэлтийг танилцууллаа. Дараагийн хэсэгт дипломын ажлын нэгдсэн дүгнэлт, цаашдын хөгжүүлэлтийн чиглэлийг авч үзнэ.
